\label{chap:planck-autodual}

\section{Compton realization associated with the temporal periodicity of \(108\) phases}

On the cycle of \cref{eq:realizacion-circular-ct108}, we reuse the positive action
section \(\hbar\) already generated by the determinantal action branch.
The associated angular quantum is
\begin{equation}
E_{108}=\hbar\omega_{108},
\qquad
m_{108}=\frac{E_{108}}{c^2}
=\frac{\hbar\omega_{108}}{c^2}.
\label{eq:cuanto-ct108}
\end{equation}
\(G\) has not yet been used. The geometry of the cycle and angular
quantization suffice to fix its two Compton readers.

\begin{theorem}[Identity between the temporal periodicity of \(108\) phases and the Compton length]
\label{thm:ct108-compton}
For the quantum \eqref{eq:cuanto-ct108},
\begin{equation}
\lambdabarC(m_{108})=\frac{\hbar}{m_{108}c}=R_{108},
\qquad
\lambdaC(m_{108})=\frac{h}{m_{108}c}=C_{108}.
\label{eq:identidad-ct108-compton}
\end{equation}
Therefore, the cycle radius is its reduced Compton wavelength, and the
length of one revolution is its full Compton wavelength.
\end{theorem}

\begin{proof}
Substituting \(m_{108}=\hbar\omega_{108}/c^2\),
\[
\frac{\hbar}{m_{108}c}
=\frac{\hbar c^2}{\hbar\omega_{108}c}
=\frac{c}{\omega_{108}}=R_{108}.
\]
Since \(h=2\pi\hbar\), the second equality is \(2\pi R_{108}=C_{108}\).
\end{proof}

The coincidence is not a numerical fit. It is a homogeneous identity: it holds
for all \(\hbar,c,t_0>0\) and depends only on the same period determining
the frequency, the cycle length, and the quantum of action.

\section{Gravitational Transducer and Self-Dual Selection Equation}

The gravitational line HMT--MD contains the transducer
\begin{equation}
G(\theta)=\theta\frac{c^5t_0^2}{\hbar},
\qquad \theta>0.
\label{eq:transductor-g-planck-anexo}
\end{equation}
Homogeneity is exact, but by itself does not select \(\theta\). The
CT108 cycle supplies an additional internal condition: the gravitational reader
must return the same radial scale as the geometric and Compton
readers.

We define the reduced gravitational radius
\begin{equation}
\rg(\theta):=\frac{G(\theta)m_{108}}{c^2}.
\label{eq:radio-gravitatorio-reducido}
\end{equation}
Substituting \eqref{eq:cuanto-ct108} and
\eqref{eq:transductor-g-planck-anexo},
\begin{equation}
\rg(\theta)
=\theta\frac{\pi}{54}\ell_0.
\label{eq:rg-theta-lineal}
\end{equation}

\begin{theorem}[Self-dual selector of the temporal periodicity of \(108\) phases]
\label{thm:selector-autodual-ct108}
The equation
\[
\rg(\theta)=R_{108}
\]
has a unique positive solution:
\begin{equation}
\boxed{\thetaStar=\left(\frac{54}{\pi}\right)^2.}
\label{eq:theta-star}
\end{equation}
\end{theorem}

\begin{proof}
By \eqref{eq:realizacion-circular-ct108} and
\eqref{eq:rg-theta-lineal}, the equation is
\[
\theta\frac{\pi}{54}\ell_0
=\frac{54}{\pi}\ell_0.
\]
Since \(\ell_0>0\), simplification yields
\(\theta=(54/\pi)^2\). The function
\(\theta\mapsto\rg(\theta)\) is strictly increasing on
\(\RR_{>0}\), so there is no other positive solution.
\end{proof}

\begin{corollary}[Internal Planck Unit]
\label{cor:unidad-planck-interna}
With \(G_*=G(\thetaStar)\), the Planck sections satisfy
\begin{equation}
\ellP:=\sqrt{\frac{\hbar G_*}{c^3}}=R_{108},
\qquad
\tP:=\sqrt{\frac{\hbar G_*}{c^5}}=\frac1{\omega_{108}},
\label{eq:planck-long-time-ct108}
\end{equation}
\begin{equation}
\mP:=\sqrt{\frac{\hbar c}{G_*}}=m_{108},
\qquad
\EP:=\mP c^2=E_{108}.
\label{eq:planck-mass-energy-ct108}
\end{equation}
In particular,
\begin{equation}
\boxed{
R_{108}=\lambdabarC=\rg=\ellP,
\qquad
C_{108}=\lambdaC=2\pi\rg=2\pi\ellP.}
\label{eq:tres-lectores-planck}
\end{equation}
\end{corollary}

\begin{proof}
\(G_*=(54/\pi)^2c^5t_0^2/\hbar\) is substituted. Then
\[
\ellP=\frac{54}{\pi}ct_0=R_{108},
\qquad
\tP=\frac{54}{\pi}t_0=\frac1{\omega_{108}},
\]
and
\[
\mP=\frac{\pi}{54}\frac{\hbar}{c^2t_0}
=\frac{\hbar\omega_{108}}{c^2}=m_{108}.
\]
The rest follows from \cref{thm:ct108-compton}.
\end{proof}

\begin{remark}[Three readers, one section]
Circular geometry publishes \(R_{108}\) and \(C_{108}\); Compton quantization
publishes \(\lambdabarC\) and \(\lambdaC\); gravitational realization
publishes \(\rg\) and \(2\pi\rg\). In \(\thetaStar\) those three pairs are
coordinates of a single section. This is the precise reason why the
Planck unit ceases to appear as an accidental dimensional combination
within this formalization.
\end{remark}

\section{Equivalence of normalizations on the self-dual identity}

The transducer \eqref{eq:transductor-g-planck-anexo} admits two ways of
fixing the elementary scale, which must not be confused.

\begin{enumerate}[label=\textup{(\roman*)}]
\item If \(\theta=1\) is fixed, then
\(\sqrt{\hbar G/c^3}=\ell_0\). The elementary transition \(ct_0\) is identified
with the Planck length of that chart.
\item If the \emph{complete orbit} CT108 is required to have Planck radius,
the condition selects \(\thetaStar=(54/\pi)^2\) and
\(R_{108}=\ellP\).
\end{enumerate}

They are not rival values of \(G\). They are two gauges of the same
homogeneous equation: one assigns the Planck unit to the elementary step; the other,
to the radius of the complete cycle. The annex uses the latter because its thesis concerns
the identity of the CT108 closure.

\section{Gravitational conventions: reduced radius and Schwarzschild radius}

The gravitational notation is fixed in an unambiguous manner:
\[
\rg=\frac{Gm}{c^2},
\qquad
\rs=\frac{2Gm}{c^2}=2\rg.
\]
The preceding theorem uses \(\rg\), not \(\rs\). In the usual convention of
Planck mass, the equality \(\hbar/(mc)=Gm/c^2\) is produced in \(m=\mP\),
whereas
\begin{equation}
\frac{\hbar}{mc}=\frac{2Gm}{c^2}
\quad\Longleftrightarrow\quad
m=\frac{\mP}{\sqrt2}.
\label{eq:cruce-schwarzschild-factor-dos}
\end{equation}
For this reason the phrase “gravitational scale” designates, in the rector thesis, the reduced
radius. The conventional horizon radius remains a distinct reader.

\section{Bifurcation of the action line}

The internal line of action possesses two oriented sections:
\begin{equation}
\hbarpre=H_5\,10^{-34}\mathcal U_S,
\qquad
\hbarret=\eta_{\mathrm{ret}}\,10^{-34}\mathcal U_S,
\qquad
\Ract=\frac{H_5}{\eta_{\mathrm{ret}}}.
\label{eq:hbar-pre-ret-anexo}
\end{equation}
The decade \(10^{-34}\) fixes the absolute scale; the quotient \(\Ract\)
transports the oriented difference. This separation makes it possible to ask whether
Planck self-duality survives in both coordinates.

\begin{theorem}[Reciprocal Covariance of Action and Gravity]
\label{thm:accion-gravedad-reciproca}
For \(a\in\{\mathrm{pre},\mathrm{ret}\}\), let
\[
G_a(\theta)=\theta\frac{c^5t_0^2}{\hbar_a},
\qquad
m_{108,a}=\frac{\hbar_a\omega_{108}}{c^2}.
\]
Then
\begin{equation}
\frac{G_a(\theta)m_{108,a}}{c^2}
=\theta\frac{\pi}{54}\ell_0
\label{eq:invariante-pre-ret-planck}
\end{equation}
is independent of \(a\). Moreover,
\begin{equation}
\frac{m_{108,\mathrm{pre}}}{m_{108,\mathrm{ret}}}=\Ract,
\qquad
\frac{G_{\mathrm{pre}}}{G_{\mathrm{ret}}}=\Ract^{-1},
\label{eq:reciprocidad-pre-ret}
\end{equation}
and the Planck length
\(\sqrt{\hbar_aG_a/c^3}=\sqrt\theta\,\ell_0\) is the same in both
orientations.
\end{theorem}

\begin{proof}
The factors \(\hbar_a\) cancel in the product \(G_am_{108,a}\).
The two ratios in \eqref{eq:reciprocidad-pre-ret} follow directly
from the definitions. The invariance of the Planck length follows from
\(\hbar_aG_a=\theta c^5t_0^2\).
\end{proof}

A symmetric parametrization makes the geometry visible:
\begin{equation}
\hbar_{\mathrm{geom}}:=\sqrt{\hbarpre\hbarret},
\qquad
\xi:=\frac12\log\Ract,
\qquad
\hbarpre=\hbar_{\mathrm{geom}}e^\xi,
\quad
\hbarret=\hbar_{\mathrm{geom}}e^{-\xi}.
\label{eq:rapidez-accion-pre-ret}
\end{equation}
Bidirectional memory is lodged in \(\xi\); Planck geometry is
lodged in the invariant product. This allows a real difference between the
two action coordinates without requiring a difference in
propagation speed.

\begin{remark}[Physical scope of the theorem]
The reciprocity of \cref{thm:accion-gravedad-reciproca} is exact within the transducer. To identify \(G_{\mathrm{pre}}\) and \(G_{\mathrm{ret}}\) with two oriented gravitational measurements, the protocol, the observable, and the symmetry relating them must be specified. The theorem proves structural compatibility; by itself, it neither anticipates a Lorentz violation nor an observable variation of \(G\).
\end{remark}
